% Options for packages loaded elsewhere
% Options for packages loaded elsewhere
\PassOptionsToPackage{unicode}{hyperref}
\PassOptionsToPackage{hyphens}{url}
\PassOptionsToPackage{dvipsnames,svgnames,x11names}{xcolor}
%
\documentclass[
  12pt,
  letterpaper]{article}
\usepackage{xcolor}
\usepackage[margin=1in]{geometry}
\usepackage{amsmath,amssymb}
\setcounter{secnumdepth}{5}
\usepackage{iftex}
\ifPDFTeX
  \usepackage[T1]{fontenc}
  \usepackage[utf8]{inputenc}
  \usepackage{textcomp} % provide euro and other symbols
\else % if luatex or xetex
  \usepackage{unicode-math} % this also loads fontspec
  \defaultfontfeatures{Scale=MatchLowercase}
  \defaultfontfeatures[\rmfamily]{Ligatures=TeX,Scale=1}
\fi
\usepackage{lmodern}
\ifPDFTeX\else
  % xetex/luatex font selection
\fi
% Use upquote if available, for straight quotes in verbatim environments
\IfFileExists{upquote.sty}{\usepackage{upquote}}{}
\IfFileExists{microtype.sty}{% use microtype if available
  \usepackage[]{microtype}
  \UseMicrotypeSet[protrusion]{basicmath} % disable protrusion for tt fonts
}{}
\usepackage{setspace}
\makeatletter
\@ifundefined{KOMAClassName}{% if non-KOMA class
  \IfFileExists{parskip.sty}{%
    \usepackage{parskip}
  }{% else
    \setlength{\parindent}{0pt}
    \setlength{\parskip}{6pt plus 2pt minus 1pt}}
}{% if KOMA class
  \KOMAoptions{parskip=half}}
\makeatother
% Make \paragraph and \subparagraph free-standing
\makeatletter
\ifx\paragraph\undefined\else
  \let\oldparagraph\paragraph
  \renewcommand{\paragraph}{
    \@ifstar
      \xxxParagraphStar
      \xxxParagraphNoStar
  }
  \newcommand{\xxxParagraphStar}[1]{\oldparagraph*{#1}\mbox{}}
  \newcommand{\xxxParagraphNoStar}[1]{\oldparagraph{#1}\mbox{}}
\fi
\ifx\subparagraph\undefined\else
  \let\oldsubparagraph\subparagraph
  \renewcommand{\subparagraph}{
    \@ifstar
      \xxxSubParagraphStar
      \xxxSubParagraphNoStar
  }
  \newcommand{\xxxSubParagraphStar}[1]{\oldsubparagraph*{#1}\mbox{}}
  \newcommand{\xxxSubParagraphNoStar}[1]{\oldsubparagraph{#1}\mbox{}}
\fi
\makeatother


\usepackage{longtable,booktabs,array}
\usepackage{calc} % for calculating minipage widths
% Correct order of tables after \paragraph or \subparagraph
\usepackage{etoolbox}
\makeatletter
\patchcmd\longtable{\par}{\if@noskipsec\mbox{}\fi\par}{}{}
\makeatother
% Allow footnotes in longtable head/foot
\IfFileExists{footnotehyper.sty}{\usepackage{footnotehyper}}{\usepackage{footnote}}
\makesavenoteenv{longtable}
\usepackage{graphicx}
\makeatletter
\newsavebox\pandoc@box
\newcommand*\pandocbounded[1]{% scales image to fit in text height/width
  \sbox\pandoc@box{#1}%
  \Gscale@div\@tempa{\textheight}{\dimexpr\ht\pandoc@box+\dp\pandoc@box\relax}%
  \Gscale@div\@tempb{\linewidth}{\wd\pandoc@box}%
  \ifdim\@tempb\p@<\@tempa\p@\let\@tempa\@tempb\fi% select the smaller of both
  \ifdim\@tempa\p@<\p@\scalebox{\@tempa}{\usebox\pandoc@box}%
  \else\usebox{\pandoc@box}%
  \fi%
}
% Set default figure placement to htbp
\def\fps@figure{htbp}
\makeatother





\setlength{\emergencystretch}{3em} % prevent overfull lines

\providecommand{\tightlist}{%
  \setlength{\itemsep}{0pt}\setlength{\parskip}{0pt}}



 


\usepackage{setspace}
\usepackage{booktabs}
\usepackage{caption}
\usepackage{longtable}
\usepackage{array}
\usepackage{multirow}
\usepackage{float}
\usepackage{graphicx}
\usepackage{threeparttable}
\usepackage{threeparttablex}
\captionsetup{width=\textwidth}
\setlength\LTcapwidth{\textwidth}
\floatplacement{figure}{H}
\floatplacement{table}{H}
\raggedright
\usepackage{setspace}
\usepackage{booktabs}
\usepackage{caption}
\usepackage{longtable}
\usepackage{array}
\usepackage{multirow}
\usepackage{float}
\usepackage{graphicx}
\usepackage{threeparttable}
\usepackage{threeparttablex}
\usepackage{siunitx}
\captionsetup{width=\textwidth}
\setlength\LTcapwidth{\textwidth}
\floatplacement{figure}{H}
\floatplacement{table}{H}
\raggedright
\makeatletter
\@ifpackageloaded{caption}{}{\usepackage{caption}}
\AtBeginDocument{%
\ifdefined\contentsname
  \renewcommand*\contentsname{Table of contents}
\else
  \newcommand\contentsname{Table of contents}
\fi
\ifdefined\listfigurename
  \renewcommand*\listfigurename{List of Figures}
\else
  \newcommand\listfigurename{List of Figures}
\fi
\ifdefined\listtablename
  \renewcommand*\listtablename{List of Tables}
\else
  \newcommand\listtablename{List of Tables}
\fi
\ifdefined\figurename
  \renewcommand*\figurename{Figure}
\else
  \newcommand\figurename{Figure}
\fi
\ifdefined\tablename
  \renewcommand*\tablename{Table}
\else
  \newcommand\tablename{Table}
\fi
}
\@ifpackageloaded{float}{}{\usepackage{float}}
\floatstyle{ruled}
\@ifundefined{c@chapter}{\newfloat{codelisting}{h}{lop}}{\newfloat{codelisting}{h}{lop}[chapter]}
\floatname{codelisting}{Listing}
\newcommand*\listoflistings{\listof{codelisting}{List of Listings}}
\makeatother
\makeatletter
\makeatother
\makeatletter
\@ifpackageloaded{caption}{}{\usepackage{caption}}
\@ifpackageloaded{subcaption}{}{\usepackage{subcaption}}
\makeatother
\usepackage{bookmark}
\IfFileExists{xurl.sty}{\usepackage{xurl}}{} % add URL line breaks if available
\urlstyle{same}
\hypersetup{
  pdftitle={Online Appendix --- The Accountability Paradox: Why South Korea's Impeachment Failed to Restore Institutional Trust},
  pdfauthor={Jeffrey Stark},
  colorlinks=true,
  linkcolor={blue},
  filecolor={Maroon},
  citecolor={blue},
  urlcolor={blue},
  pdfcreator={LaTeX via pandoc}}


\title{Online Appendix --- The Accountability Paradox: Why South Korea's
Impeachment Failed to Restore Institutional Trust}
\author{Jeffrey Stark}
\date{March 3, 2026}
\begin{document}
\maketitle


\setstretch{2}
\section{Appendix A: Data and Measurement}\label{sec-data}

\subsection{Survey Instruments}\label{sec-instruments}

The analysis draws on two nationally representative surveys of South
Korean adults. This appendix section documents fielding details,
sampling procedures, and relevant methodological features of each
instrument.

\textbf{Asian Barometer Survey (ABS).} The ABS is a cross-national
survey of democratic values and political attitudes conducted across
eighteen Asian polities. Six waves cover South Korea: Wave 1 (2003),
Wave 2 (2006), Wave 3 (2010--2011), Wave 4 (2014--2015), Wave 5
(2015--2016), and Wave 6 (2021--2022). Each Korean wave uses multistage
stratified probability sampling with urban-rural stratification and
proportional allocation across administrative regions, yielding sample
sizes of approximately 1,200 respondents per wave. Face-to-face
interviews were conducted by trained interviewers using structured
questionnaires. The ABS is the primary source for establishing long-run
trends in democratic satisfaction, quality assessment, governmental
responsiveness, and the output legitimacy mechanism (Section 4.1 of the
main text).

\textbf{Korean Academic Multimode Open Survey (KAMOS).} KAMOS is a
mixed-mode nationally representative panel survey of South Korean
adults, combining online panel and offline probability samples to
achieve demographic representativeness. Wave 1 was fielded in
November--December 2016, during the active protest period preceding
President Park Geun-hye's National Assembly impeachment vote (December
9, 2016). Wave 4 was fielded in 2019, approximately two years into the
Moon Jae-in administration. The institutional trust battery, fielded in
both waves, rates ten institutions on a five-point scale (1 = ``do not
trust at all,'' 5 = ``trust very much''). KAMOS is the primary source
for the trust architecture analysis and all regression models reported
in the main text.

Table A1 presents wave-by-wave sample sizes and fielding windows for
both surveys.

\subsection{KAMOS Institutional Trust Battery: Full Question
Wording}\label{sec-question-wording}

The KAMOS institutional trust battery asks respondents: \emph{``How much
do you trust each of the following institutions?''} Responses are
recorded on a five-point scale: 1 = Do not trust at all, 2 = Do not
trust, 3 = Neutral, 4 = Trust, 5 = Trust very much.

The ten institutions and their functional classification for this
analysis are as follows. Executive institutions (Category A): central
government (\emph{중앙정부}); local government (\emph{지방정부}).
Horizontal accountability institutions (Category B): National Assembly
(\emph{국회}); legislature (\emph{입법부}, used interchangeably in some
KAMOS waves). Societal accountability institutions (Category C): media
(\emph{언론}); NGOs/civil society organizations (\emph{시민단체}).
Excluded from primary analysis (Category D): private enterprise
(\emph{기업}); religious institutions (\emph{종교기관}). Retained for
robustness checks: trust in society generally (\emph{사회 전반}); trust
in fellow citizens (\emph{타인}).

The Category B composite (National Assembly + legislature) is computed
as the simple mean of the two items when both are available; Cronbach's
α = 0.78, Pearson \emph{r} = 0.64 across the pooled sample.

\subsection{Institution-Type Classification
Rationale}\label{sec-classification}

The three-category classification used in the main analysis reflects
functional distinctions within the democratic accountability
architecture rather than arbitrary grouping. Executive institutions
(Category A) are those whose primary function is policy implementation
and whose evaluative standard --- in both normative democratic theory
and, the argument in the main text claims, in Korean political culture
--- is performance output. Horizontal accountability institutions
(Category B) are those whose constitutional function is to constrain and
oversee the executive; their evaluative standard is procedural --- the
quality of deliberation, the independence of adjudication, the diligence
of oversight. Societal accountability institutions (Category C) perform
information provision and interest articulation functions that are
prerequisites for meaningful horizontal accountability; their
legitimacy, like that of Category B institutions, rests on procedural
rather than output criteria.

Category D institutions (private enterprise, religious) are excluded
because their relationship to the democratic accountability architecture
is indirect at best. Private enterprise trust is more plausibly governed
by economic conditions and personal financial experience than by
democratic institutional performance; religious trust reflects
denominational affiliation and spiritual community membership rather
than evaluations of political authority. Including these items in the
primary analysis would contaminate the institution-type contrast with
theoretically unrelated variation.

\subsection{Full KAMOS Variable Direction
Table}\label{sec-direction-table}

Table A2 presents wave-level means for all KAMOS variables used in or
relevant to this analysis, with direction-of-change indicators. The
table documents the central descriptive pattern motivating the
accountability paradox argument: the simultaneous rise in
output-oriented indicators (economic evaluations, political
satisfaction, national pride, electoral engagement) and fall in
procedural institutional trust.

\subsection{Sample Descriptive Statistics}\label{sec-descriptives}

Table A3 presents descriptive statistics for all variables entering the
regression analyses, pooled across KAMOS Waves 1 and 4. Variables have
been rescaled as described in Section 3.2 of the main text prior to
model estimation.

\section{Appendix B: Extended Main Results}\label{sec-main-results}

\subsection{Full Three-Way Trust Model
Coefficients}\label{sec-full-threeway}

Table 3 in the main text reports the key interaction terms from the
three-way pooled OLS model but suppresses coefficients for control
variables in the interest of presentation. Table B1 here presents the
full coefficient table for the same model, including main effects for
all control variables.

\subsection{Output Legitimacy Moderation: Full
Results}\label{sec-output-legitimacy-full}

Table B2 presents full results from the output legitimacy moderation
model, which extends the three-way specification by adding interactions
between economic evaluation and institution type. This model provides
the primary test of the second component of H1: that economic optimism
selectively buffers executive institutional trust while providing no
equivalent protection for intermediary institutions.

The economic evaluation composite is the mean of the family economic
conditions and national economic conditions items, normalized to a 0--1
scale. All items are centered at their grand mean prior to interaction
construction to improve interpretability of lower-order coefficients.
Three models are reported: Model 1 uses family economic conditions
alone; Model 2 uses national economic conditions alone; Model 3 uses the
composite. Consistency across operationalizations is expected under the
output legitimacy account, since both items tap the same underlying
orientation.

\begin{table}

\caption{\label{tbl-output-legitimacy-full}Output legitimacy moderation
model --- full coefficient table. Dependent variable: institutional
trust (0--1 scale). Three-way interaction model with economic evaluation
composite. Standard errors clustered by respondent. Source: KAMOS Waves
1 and 4.}

\centering{

\centering
\begin{tabular}[t]{lc}
\toprule
  & (1)\\
\midrule
Wave 4 (2019) & \num{-0.0228}\\
 & (\num{0.0220})\\
Horizontal institution & \num{-0.1051}***\\
 & (\num{0.0127})\\
Societal institution & \num{-0.0293}*\\
 & (\num{0.0118})\\
Economic evaluation & \num{-0.1702}***\\
 & (\num{0.0208})\\
Wave 4 × Horizontal & \num{-0.0652}**\\
 & (\num{0.0203})\\
Wave 4 × Societal & \num{-0.0150}\\
 & (\num{0.0191})\\
Wave 4 × Econ (Executive shield) & \num{0.0219}\\
 & (\num{0.0309})\\
Horizontal × Econ & \num{-0.0154}\\
 & (\num{0.0181})\\
Societal × Econ & \num{0.0877}***\\
 & \vphantom{1} (\num{0.0174})\\
Wave 4 × Econ × Horizontal & \num{0.0330}\\
 & (\num{0.0272})\\
Wave 4 × Econ × Societal & \num{-0.0599}*\\
 & (\num{0.0261})\\
Age & \num{0.0297}\\
 & (\num{0.0174})\\
Gender (male = 1) & \num{-0.0053}\\
 & (\num{0.0052})\\
Education & \num{-0.0416}*\\
 & (\num{0.0207})\\
Income & \num{0.0447}**\\
 & (\num{0.0164})\\
Constant & \num{0.6054}***\\
 & (\num{0.0210})\\
\midrule
Num.Obs. & \num{17495}\\
R2 & \num{0.117}\\
R2 Adj. & \num{0.116}\\
\bottomrule
\multicolumn{2}{l}{\rule{0pt}{1em}* p $<$ 0.05, ** p $<$ 0.01, *** p $<$ 0.001}\\
\multicolumn{2}{l}{\rule{0pt}{1em}OLS with survey weights; SEs clustered by respondent (HC1). Reference: Executive institution, Wave 1 (2016). Economic evaluation = composite of family and national economic assessments (0–1). * p $<$ .05, ** p $<$ .01, *** p $<$ .001.}\\
\end{tabular}

}

\end{table}%

The key quantities of interest are the Economic Evaluation × Wave ×
Institution Type three-way interactions, which test whether the change
in the relationship between economic evaluation and trust across waves
differed by institution type. A positive and significant Economic
Evaluation × Wave × Executive coefficient would indicate that economic
optimism buffered executive trust specifically over the 2016--2019
period. Null or negative analogous coefficients for the intermediary
categories would confirm the category-specificity of the buffer --- that
is, confirm the output legitimacy mechanism rather than a general
positivity effect that operates uniformly across all institution types.

A contrast test of the executive buffering coefficient against each
intermediary coefficient is reported below the model table, along with
the associated \emph{t}-statistic and two-tailed \emph{p}-value.

\subsection{Partisan Moderation: Full Results}\label{sec-partisan-full}

Table B3 presents full results from the partisan moderation check. The
model adds Wave × Party Identification interactions to the three-way
specification to evaluate whether the trust collapse was concentrated
among partisan losers --- conservative supporters of President Park
Geun-hye who may have registered dissatisfaction with the accountability
institutions that removed her.

Party identification is measured on a seven-point scale ranging from
strong progressive to strong conservative, normalized to 0--1. Two
specifications are reported. Model 1 uses ideological self-placement as
a continuous moderator. Model 2 uses a categorical indicator
distinguishing strong conservatives (top quartile of the ideology
distribution) from the remainder of the sample, providing a direct test
of whether the collapse was driven by the subset of respondents most
likely to have experienced the impeachment as a partisan defeat.

\section{Appendix C: Robustness Checks}\label{sec-robustness}

\subsection{Ordered Logit Specification}\label{sec-ologit}

The OLS specification in the main analysis imposes the assumption of
equal intervals between trust scale points --- a convenient
approximation that may not hold for ordinal Likert items. Table C1
reports results from an ordered logit model that respects the ordinal
structure of the five-point trust scale without assuming equal spacing.

Ordered logit coefficients are not directly comparable to OLS estimates
in magnitude, but the pattern of sign, significance, and relative
magnitude across institution-type interactions is directly
interpretable. The key test is whether the Wave × Horizontal
Accountability and Wave × Societal Accountability interaction
coefficients remain significantly negative relative to the executive
reference category, and whether their magnitude remains substantially
larger than the baseline Wave coefficient.

\subsection{Alternative Institution-Type
Groupings}\label{sec-alt-groupings}

The primary analysis assigns ten KAMOS trust items to four categories,
using six items in the main models (Categories A--C) and reserving two
(trust in society, trust in fellow citizens) as robustness checks. The
exclusion of the latter two items is theoretically motivated --- they
plausibly capture generalized social trust rather than political
institutional confidence --- but is a classification decision whose
sensitivity warrants examination.

Table C2 reports three alternative grouping specifications. Model 1
includes trust in society and trust in fellow citizens as a fourth
institution-type category (Diffuse Social Trust), added to the three-way
model with its own Wave interaction. Model 2 reassigns trust in society
to the societal accountability category (Category C), treating it as an
aggregate indicator of the associational environment in which civil
society operates. Model 3 applies the most permissive grouping, treating
all ten items as separate institution types with separate Wave
interactions, estimating ten institution-specific change coefficients
relative to the executive reference.

The key finding to confirm is that the executive--intermediary
divergence pattern documented in the main analysis is not an artifact of
the specific grouping decisions made there.

\begin{table}

\caption{\label{tbl-alt-groupings}Alternative institution-type grouping
specifications. Model 1 (baseline): original three-category
specification. Model 2: trust in society added to the Societal category.
Model 3: trust in society treated as a separate fourth category.
Executive institutions are the reference category in all models.
Standard errors clustered by respondent. Source: KAMOS Waves 1 and 4.}

\centering{

\centering
\begin{tabular}[t]{lccc}
\toprule
  & Model 1 (Baseline) & Model 2 (Society→Cat. C) & Model 3 (Society=Cat. D)\\
\midrule
Wave 4 (2019) & \num{-0.0206}*** & \num{-0.0208}*** & \num{-0.0208}***\\
 & (\num{0.0060}) & (\num{0.0060}) & (\num{0.0060})\\
Horizontal institution & \num{-0.1151}*** & \num{-0.1151}*** & \num{-0.1151}***\\
 & (\num{0.0035}) & (\num{0.0035}) & (\num{0.0035})\\
Societal institution & \num{0.0277}*** & \num{0.0362}*** & \num{0.0277}***\\
 & (\num{0.0032}) & (\num{0.0028}) & (\num{0.0032})\\
Social trust (added item) &  &  & \num{0.0532}***\\
 &  &  & (\num{0.0036})\\
Wave 4 × Horizontal & \num{-0.0424}*** & \num{-0.0424}*** & \num{-0.0424}***\\
 & (\num{0.0055}) & (\num{0.0055}) & (\num{0.0055})\\
Wave 4 × Societal & \num{-0.0516}*** & \num{-0.0333}*** & \num{-0.0516}***\\
 & (\num{0.0049}) & (\num{0.0043}) & (\num{0.0049})\\
Wave 4 × Social trust &  &  & \num{0.0033}\\
 &  &  & (\num{0.0055})\\
Constant & \num{0.4977}*** & \num{0.4937}*** & \num{0.4937}***\\
 & (\num{0.0176}) & (\num{0.0168}) & (\num{0.0168})\\
\midrule
Num.Obs. & \num{17495} & \num{20994} & \num{20994}\\
R2 & \num{0.096} & \num{0.096} & \num{0.106}\\
\bottomrule
\multicolumn{4}{l}{\rule{0pt}{1em}* p $<$ 0.05, ** p $<$ 0.01, *** p $<$ 0.001}\\
\multicolumn{4}{l}{\rule{0pt}{1em}OLS with survey weights; SEs clustered by respondent (HC1). Reference: Executive institution, Wave 1 (2016). Controls (age, gender, education, income, ideology) included but not shown. * p $<$ .05, ** p $<$ .01, *** p $<$ .001.}\\
\end{tabular}

}

\end{table}%

\subsection{Category D Placebo Test}\label{sec-placebo}

A key identification concern for the output legitimacy account is that
the economic evaluation moderation might reflect a general economic
positivity halo --- citizens who evaluate the economy favorably simply
trust all institutions more, rather than executive institutions
specifically. If so, the ``executive shield'' is a measurement artifact
rather than evidence of category-specific output legitimacy. Table C3
reports the placebo test using Category D institutions (private
enterprise, religious institutions), whose trust trajectory should be
governed by economic conditions and personal experience rather than
democratic performance frames. Under the output legitimacy account, the
Wave × Economic Evaluation interaction should be positive for executive
institutions but null for Category D institutions --- just as it should
be null for horizontal and societal accountability institutions.

\begin{table}

\caption{\label{tbl-placebo}Category D placebo test. Output legitimacy
moderation model extended to include private enterprise and religious
institutions (Category D) as a fourth institution type. Dependent
variable: institutional trust (0--1 scale). The `Wave 4 × Econ ×
Category D' row provides the placebo test of the executive shield
hypothesis. Standard errors clustered by respondent. Source: KAMOS Waves
1 and 4.}

\centering{

\centering
\begin{tabular}[t]{lc}
\toprule
  & (1)\\
\midrule
Wave 4 (2019) & \num{-0.0221}\\
 & (\num{0.0220})\\
Horizontal institution & \num{-0.1051}***\\
 & (\num{0.0127})\\
Societal institution & \num{-0.0293}*\\
 & (\num{0.0118})\\
Category D institution & \num{-0.0103}\\
 & (\num{0.0123})\\
Economic evaluation & \num{-0.1691}***\\
 & (\num{0.0208})\\
Wave 4 × Econ (Executive shield) & \num{0.0195}\\
 & (\num{0.0309})\\
Wave 4 × Horizontal & \num{-0.0652}**\\
 & \vphantom{1} (\num{0.0203})\\
Wave 4 × Societal & \num{-0.0150}\\
 & (\num{0.0191})\\
Wave 4 × Category D & \num{-0.0539}**\\
 & (\num{0.0196})\\
Wave 4 × Econ × Horizontal & \num{0.0330}\\
 & (\num{0.0272})\\
Wave 4 × Econ × Societal & \num{-0.0599}*\\
 & (\num{0.0261})\\
Wave 4 × Econ × Category D (placebo) & \num{0.0168}\\
 & (\num{0.0269})\\
Constant & \num{0.5937}***\\
 & (\num{0.0203})\\
\midrule
Num.Obs. & \num{24493}\\
R2 & \num{0.105}\\
\bottomrule
\multicolumn{2}{l}{\rule{0pt}{1em}* p $<$ 0.05, ** p $<$ 0.01, *** p $<$ 0.001}\\
\multicolumn{2}{l}{\rule{0pt}{1em}OLS with survey weights; SEs clustered by respondent (HC1). Reference: Executive institution, Wave 1 (2016). Category D = private enterprise + religious institutions. Placebo prediction: Wave 4 × Econ × Category D should be null or negative. Controls (age, gender, education, income) included but not shown. * p $<$ .05, ** p $<$ .01, *** p $<$ .001.}\\
\end{tabular}

}

\end{table}%

\subsection{Alternative Economic Evaluation
Operationalizations}\label{sec-econ-alt}

The output legitimacy moderation analysis uses a composite of family and
national economic evaluation. Two concerns motivate additional checks.
First, the composite weights the two items equally; it is possible that
the output legitimacy mechanism operates through one item but not the
other, which equal weighting would obscure. Second, political
satisfaction --- which also rose between 2016 and 2019 --- is
conceptually adjacent to output legitimacy but not identical to it; a
citizen could be satisfied with the Moon administration for reasons
other than economic performance (e.g., satisfaction with the
accountability outcome itself).

Table C3 reports the output legitimacy moderation model under three
alternative economic operationalizations. Model 1 uses family economic
conditions alone (personal pocketbook assessment). Model 2 uses national
economic conditions alone (sociotropic assessment). Model 3 uses
political satisfaction as the moderator rather than economic evaluation,
providing a test of whether the moderation reflects output legitimacy
specifically or a more general positivity orientation. Under the output
legitimacy account, Models 1 and 2 should replicate the main pattern;
Model 3 is expected to attenuate it, since political satisfaction in the
post-impeachment period likely reflects satisfaction with the
accountability outcome as much as economic contentment.

\subsection{ABS Wave 5 Dip and Wave 6 Partial
Recovery}\label{sec-abs-robustness}

The KAMOS data, covering 2016 and 2019 only, cannot resolve whether the
trust decline within the ABS institutional trust battery was
concentrated in or around the impeachment period or represents a more
gradual erosion. The ABS battery provides independent evidence on this
question: Wave 5, whose Korean fieldwork occurred in 2015--2016
overlapping with the beginning of the protest period, shows a trust dip
across multiple institutional categories; Wave 6, fielded in 2021--2022
under a different political context, shows partial but incomplete
recovery.

Figure C1 plots ABS wave-level means for the full institutional trust
battery across all six waves. The figure serves two purposes. First, it
provides independent confirmation --- from a different survey
instrument, with different items and different fieldwork timing --- that
the impeachment period coincided with a trust shock of the type
documented in KAMOS. Second, the incomplete Wave 6 recovery speaks to
the durability of the pattern: a partial rebound after 2019 would be
consistent with the non-recovery account (not a full restoration of
pre-crisis trust) and inconsistent with a simple temporary-shock account
(which would predict full recovery once political conditions
normalized).

\subsection{Excluding the Transitional Wave
Window}\label{sec-wave-window}

A potential concern with the KAMOS design is that Wave 1
(November--December 2016) was fielded during an active political crisis
--- the protest movement was ongoing, the National Assembly had not yet
voted on impeachment, and political uncertainty was at a maximum. Trust
ratings in Wave 1 may therefore reflect crisis-period emotional valence
rather than a stable baseline, which would make the Wave 1 to Wave 4
comparison a contrast between crisis-period trust and post-crisis trust
rather than a pre-crisis to post-crisis comparison.

This concern is difficult to fully resolve with the available data.
However, two observations suggest it does not drive the main findings.
First, the ABS Wave 5 data --- fielded slightly earlier than KAMOS Wave
1 but covering the overlapping period --- show that Korean institutional
trust was already declining prior to the peak of the protest movement,
indicating that the KAMOS Wave 1 baseline is not artificially depressed
by crisis-moment measurement. Second, if crisis-period emotional
negativity inflated the Wave 1 to Wave 4 decline by depressing the
baseline, the pattern should be most pronounced for executive
institutions --- the primary target of citizen anger during the protests
--- rather than for intermediary institutions. The observed pattern runs
in the opposite direction, which is inconsistent with a crisis-baseline
artifact explanation.

\subsection{Weighted versus Unweighted Comparison}\label{sec-weights}

The main analysis uses KAMOS survey weights to account for the
mixed-mode design's differential coverage across demographic groups.
Table C4 compares weighted and unweighted estimates for the three-way
trust model's key interaction coefficients. Substantial divergence
between the two would indicate that demographic composition differences
across modes or across waves are driving the results rather than genuine
attitudinal change.

\section{Appendix D: Additional Descriptive
Evidence}\label{sec-descriptive-additional}

\subsection{ABS Long-Run Trends: Extended Variable
Set}\label{sec-abs-extended}

Table D1 extends Table 1 from the main text to include the full set of
ABS variables relevant to the output legitimacy mechanism and the
satisfaction-quality decoupling. In addition to democratic satisfaction,
quality assessment, and economic evaluation, the table reports
wave-level means for governmental responsiveness (\emph{Does the
government respond to what people want?}), information withholding
(\emph{The government often withholds information}), and trust in
specific institutions from the ABS battery (parliament, courts,
political parties, civil service, national government). All items are
normalized to a 0--1 scale for cross-item comparability.

\subsection{KAMOS Trust Means: Full Ten-Item
Table}\label{sec-kamos-full}

Table D2 presents wave-level means for all ten KAMOS trust items
individually, rather than the category-level aggregates reported in
Table 2 of the main text. This allows inspection of within-category
heterogeneity --- whether, for example, trust in local government moved
differently from trust in central government within the executive
category, or whether media trust diverged from NGO trust within the
societal accountability category.

\subsection{Rising KAMOS Indicators, 2016--2019}\label{sec-rising}

Table D3 documents the KAMOS indicators that rose between Wave 1 and
Wave 4 --- the context within which the trust collapse occurred. The
rising indicators include economic evaluations (family and national
conditions), political satisfaction, social mobility optimism, national
pride, and party identification. The simultaneous rise in
output-oriented evaluations and fall in procedural institutional trust
is the core empirical pattern this article explains.

\section*{References}\label{references}
\addcontentsline{toc}{section}{References}

\phantomsection\label{refs}




\end{document}
